%% RESUMO NA LINGUA VERNACULA (elemento obrigatorio -- NBR 6028)
%%
%% Regras:
%%   - paragrafo unico, sem recuo, justificado;
%%   - Times New Roman 12, espacamento simples;
%%   - entre 150 e 500 palavras;
%%   - verbo na voz ativa e na terceira pessoa do singular;
%%   - as palavras-chave sao definidas em config/dados.tex.
%%
%% Um bom resumo informa: objetivo geral; principais conceitos e autores da
%% fundamentacao teorica; etapas e acoes da execucao; populacao envolvida e
%% periodo; principais resultados alcancados; conclusoes.

\begin{resumo}
Esta dissertação propõe e aplica um modelo multidimensional para avaliar
estratégias de contextualização documental em sistemas de geração aumentada por
recuperação (RAG) usados na consulta a documentos institucionais públicos
brasileiros. O problema de pesquisa é que os órgãos públicos escolhem entre
estratégias como \textit{chunking} fixo, recursivo, semântico e
\textit{parent-child} sensível à estrutura sem um instrumento que compare, ao
mesmo tempo, qualidade da resposta, custo, rastreabilidade da fonte e risco
institucional. A pesquisa pergunta, portanto, como propor e aplicar um modelo
multidimensional que avalie essas estratégias nesses documentos. O modelo
reúne cinco dimensões --- recuperação, geração, eficiência, rastreabilidade e
adequação institucional --- e tem como contribuição principal o Índice de
Rastreabilidade Documental (IRD), composto por seis critérios binários e
validado por concordância entre avaliadores independentes e por validade
convergente com especialistas de órgãos públicos. A aplicação ocorre em um
experimento pareado com quatro estratégias, uma ablação de metadados e um
conjunto de 60 a 100 perguntas de referência, com o modelo gerador mantido
constante. Como produto técnico, a pesquisa entrega uma matriz de decisão para
o comitê de governança de dados e de inteligência artificial do órgão, avaliada
quanto à eficácia e de forma formativa com representantes dos perfis usuários.
\end{resumo}
